\chapter{Complexity — Mult}
%%% generated by the blueprint pipeline from TCSlib.Complexity.SpaceComplexity.Logspace.Mult
\section{Overview}

Declarations of the Lean module \texttt{TCSlib.Complexity.SpaceComplexity.Logspace.Mult}.

\section{Declarations}

\begin{definition}[Multiplication language]
\lean{Complexity.multLang}
\label{Complexity.multLang}
\leanok
\uses{Turing.pairEncode}
A Boolean string belongs to the multiplication language precisely when there are natural
numbers $a$ and $b$ such that it is the paired encoding of the binary digit list of $a$
and the paired encoding of the binary digit lists of $b$ and $ab$.
\end{definition}

\begin{definition}[Multiplication column sum]
\lean{Complexity.multColumn}
\label{Complexity.multColumn}
\leanok
For Boolean lists $a,b$ and a natural-number column index $j$, define the column sum to
be the sum over positions $p$ in the range from $0$ to $|a|-1$ of $a_p b_{j-p}$ when
$p\le j$, and $0$ otherwise; list entries beyond either list's length are interpreted as
false, and each Boolean is identified with $0$ or $1$.
\end{definition}

\begin{lemma}[Empty multiplication column]
\lean{Complexity.multColumn_nil}
\label{Complexity.multColumn_nil}
\leanok
\uses{Complexity.multColumn}
\difficulty{1}
Let $b$ be a finite bit string and let $j$ be a nonnegative integer. The column sum
formed from the empty bit string and $b$ at position $j$ is $0$.
\end{lemma}

\begin{lemma}[Multiplication column for a cons list]
\lean{Complexity.multColumn_cons}
\label{Complexity.multColumn_cons}
\leanok
\uses{Complexity.multColumn}
\difficulty{3}
For a bit $t$, bit strings $a,b$, and $j\in\mathbb{N}$, the multiplication-column sum
for the list beginning with $t$ and continuing with $a$ equals $t$ times the $j$th bit
of $b$, plus zero when $j=0$, and otherwise the multiplication-column sum for $a,b$ at
index $j-1$. Missing list entries are treated as zero.
\end{lemma}

\begin{lemma}[Bound on a multiplication column]
\lean{Complexity.multColumn_le}
\label{Complexity.multColumn_le}
\leanok
\uses{Complexity.multColumn}
\difficulty{4}
Let $a$ and $b$ be finite bit strings, and let $j$ be a nonnegative integer. Define the
column sum as the sum over positions $p$ in $a$ of the product of the bit of $a$ at $p$
and the bit of $b$ at position $j-p$, taking a missing bit to be $0$ and including a
term only when $p\le j$. This sum is at most the length of $a$.
\end{lemma}

\begin{lemma}[Value of a finite bit prefix]
\lean{Complexity.mult_bits_prefix}
\label{Complexity.mult_bits_prefix}
\leanok
\uses{Complexity.LogProg.bitsVal}
\difficulty{4}
For every bit string $w$ and natural number $n$, the sum of $2^j$ times the bit of $w$
at position $j$ over $0\le j<n$ equals the value of the prefix of $w$ of length $n$,
interpreting missing bits as zero.
\end{lemma}

\begin{lemma}[Leading-bit recurrence for multiplication columns]
\lean{Complexity.mult_column_prefix_cons}
\label{Complexity.mult_column_prefix_cons}
\leanok
\uses{Complexity.LogProg.bitsVal, Complexity.multColumn, Complexity.multColumn_cons, Complexity.mult_bits_prefix}
\difficulty{5}
Let $t$ be a bit, let $a,b$ be finite bit strings, and let $n\ge 0$. If
$C_j(a,b)=\sum_{p=0}^{|a|-1} [p\le j]a_p b_{j-p}$, where out-of-range bits are zero, and
$V(w)=\sum_{i=0}^{|w|-1}2^i w_i$, then
\[\sum_{j=0}^{n}2^j C_j(t::a,b)=t\,V(b\mathbin{\upharpoonright}(n+1))+2\sum_{j=0}^{n-1}2^j C_j(a,b).\]
Here $b\mathbin{\upharpoonright}(n+1)$ denotes the prefix of $b$ of length at most
$n+1$.
\end{lemma}

\begin{lemma}[Weighted multiplication columns give the product]
\lean{Complexity.mult_columns_value}
\label{Complexity.mult_columns_value}
\leanok
\uses{Complexity.LogProg.bitsVal, Complexity.multColumn, Complexity.multColumn_nil, Complexity.mult_column_prefix_cons}
\difficulty{4}
Let $a,b$ be finite bit strings and let $n$ be a natural number with $|a|+|b|\le n$. If
$C_j(a,b)=\sum_{p=0}^{|a|-1}[p\le j]a_p b_{j-p}$, with out-of-range bits taken to be
zero, then
\[\sum_{j=0}^{n-1}2^j C_j(a,b)=V(a)V(b),\]
where $V(w)=\sum_{i=0}^{|w|-1}2^i w_i$.
\end{lemma}

\begin{definition}[Multiplication carry sequence]
\lean{Complexity.multCarry}
\label{Complexity.multCarry}
\leanok
\uses{Complexity.multColumn}
For Boolean lists $a,b$, define a sequence of natural-number carries by setting the
carry at index $0$ to $0$ and, for each $j\ge 0$, setting the carry at index $j+1$ to
the quotient obtained by dividing the sum of the carry at $j$ and the multiplication
column sum at $j$ by $2$.
\end{definition}

\begin{lemma}[Bound on the multiplication carry]
\lean{Complexity.multCarry_le}
\label{Complexity.multCarry_le}
\leanok
\uses{Complexity.multCarry, Complexity.multColumn_le}
\difficulty{4}
For bit strings $a,b$ and $j\in\mathbb{N}$, the carry into column $j$ in the
multiplication procedure is at most the length of $a$.
\end{lemma}

\begin{lemma}[Invariant for multiplication carries]
\lean{Complexity.mult_carry_invariant}
\label{Complexity.mult_carry_invariant}
\leanok
\uses{Complexity.multCarry, Complexity.multColumn}
\difficulty{4}
For bit strings $a,b$ and every natural number $n$, the weighted sum of the first $n$
product-column values equals the weighted sum of the first $n$ output bits determined by
the carries, plus $2^n$ times the carry after column $n$.
\end{lemma}

\begin{lemma}[Uniqueness of binary digits]
\lean{Complexity.mult_digits_unique}
\label{Complexity.mult_digits_unique}
\leanok
\difficulty{6}
Let $n,k$ be natural numbers and let $f,g:\mathbb N\to\mathbb N$ satisfy $f(j)<2$ and
$g(j)<2$ for every $j<n$. If
\[\sum_{j=0}^{n-1}2^j f(j)+2^n k=\sum_{j=0}^{n-1}2^j g(j),\]
then $k=0$ and $f(j)=g(j)$ for every $j<n$.
\end{lemma}

\begin{lemma}[Correctness of binary multiplication columns]
\lean{Complexity.mult_columns_correct}
\label{Complexity.mult_columns_correct}
\leanok
\uses{Complexity.LogProg.bitsVal, Complexity.multCarry, Complexity.multColumn, Complexity.mult_bits_prefix, Complexity.mult_carry_invariant, Complexity.mult_columns_value, Complexity.mult_digits_unique}
\difficulty{5}
Let $a,b,c$ be finite bit strings with $|c|\le |a|+|b|$. Define the column sum
$C_j(a,b)=\sum_{p=0}^{|a|-1}[p\le j]a_p b_{j-p}$ and carries by $K_0=0$ and
$K_{j+1}=\lfloor(K_j+C_j(a,b))/2\rfloor$, taking out-of-range bits to be zero. Then
$V(c)=V(a)V(b)$ if and only if $K_{|a|+|b|}=0$ and, for every $j<|a|+|b|$, the $j$th bit
of $c$ equals $(K_j+C_j(a,b))$ modulo $2$.
\end{lemma}

\begin{lemma}[Binary length of a product]
\lean{Complexity.mult_product_length}
\label{Complexity.mult_product_length}
\leanok
\uses{Complexity.LogProg.length_bits_le}
\difficulty{3}
For all natural numbers $a$ and $b$, the length of the binary word representing $ab$ is
at most the sum of the lengths of the binary words representing $a$ and $b$.
\end{lemma}

\begin{lemma}[A single multiplication column]
\lean{Complexity.mult_column_row}
\label{Complexity.mult_column_row}
\leanok
\difficulty{4}
Let $t$ be a bit, let $b$ be a finite bit string, and let $j,p\ge 0$. Then
\[\sum_{q=0}^{|b|-1}[p+q=j]t b_q=\begin{cases}t b_{j-p},&p\le j,\0,&p>j,\end{cases}\]
where bits outside the string are taken to be zero.
\end{lemma}

\begin{lemma}[Two forms of the product-column sum]
\lean{Complexity.mult_column_double}
\label{Complexity.mult_column_double}
\leanok
\uses{Complexity.multColumn, Complexity.mult_column_row}
\difficulty{3}
For bit strings $a,b$ and every natural number $j$, summing the products of their bits
over pairs of positions $p<|a|$ and $q<|b|$ with $p+q=j$ gives the product-column value
at position $j$; missing bits are interpreted as zero.
\end{lemma}

\begin{definition}[Multiplication row sum]
\lean{Complexity.multRow}
\label{Complexity.multRow}
\leanok
For a Boolean $t$, a Boolean list $b$, and natural numbers $j,k$, define the row sum as
the sum over $q$ from $0$ to $|b|-1$ of $t b_q$ when $k+q=j$, and $0$ otherwise; list
entries beyond the list's length are interpreted as false, and Booleans as $0$ or $1$.
\end{definition}

\begin{lemma}[First-row decomposition]
\lean{Complexity.multRow_cons}
\label{Complexity.multRow_cons}
\leanok
\uses{Complexity.multRow}
\difficulty{3}
Let $t,u$ be bits, let $b$ be a finite bit string, and let $j,k$ be nonnegative
integers. Define the row sum for a bit $t$, string $b$, and indices $j,k$ by summing
over positions $q$ in $b$ the product of $t$ and the bit of $b$ at $q$ when $k+q=j$, and
$0$ otherwise; a missing bit is taken to be $0$. The row sum for $t$ and the string
formed by placing $u$ before $b$, at indices $j,k$, equals the term $t u$ when $k=j$
(and $0$ otherwise), plus the row sum for $t,b$ at indices $j,k+1$.
\end{lemma}

\begin{lemma}[Bound on a multiplication row]
\lean{Complexity.multRow_le_one}
\label{Complexity.multRow_le_one}
\leanok
\uses{Complexity.multRow, Complexity.mult_column_row}
\difficulty{4}
Let $t$ be a bit, let $b$ be a finite bit string, and let $j,k$ be nonnegative integers.
The row sum obtained by summing the product of $t$ and each bit of $b$ at position $q$
for which $k+q=j$, with all other terms zero, is at most $1$.
\end{lemma}

\begin{definition}[Sum of shifted multiplication rows]
\lean{Complexity.multRows}
\label{Complexity.multRows}
\leanok
\uses{Complexity.multRow}
For Boolean lists $a,b$ and natural numbers $j,p$, define the sum of shifted rows by
summing, over $q$ from $0$ to $|a|-1$, the row sum with Boolean entry $a_q$, list $b$,
column index $j$, and shift $p+q$; entries of $a$ beyond its length are interpreted as
false.
\end{definition}

\begin{lemma}[First-row decomposition of the row sum]
\lean{Complexity.multRows_cons}
\label{Complexity.multRows_cons}
\leanok
\uses{Complexity.multRow, Complexity.multRows}
\difficulty{3}
Let $t$ be a bit, let $a,b$ be finite bit strings, and let $j,p$ be nonnegative
integers. Define the sum of rows by summing, over positions $q$ in the first string, the
row sum formed from its bit at $q$ and the second string, at indices $j,p+q$, treating a
missing bit as $0$. The sum of rows for the string formed by placing $t$ before $a$
equals the row sum for $t,b$ at indices $j,p$, plus the sum of rows for $a,b$ at indices
$j,p+1$.
\end{lemma}

\begin{lemma}[Bound on the sum of rows]
\lean{Complexity.multRows_bound}
\label{Complexity.multRows_bound}
\leanok
\uses{Complexity.multRow_le_one, Complexity.multRows}
\difficulty{3}
Let $a,b$ be finite bit strings and let $j,p$ be nonnegative integers. For each position
$q$ in $a$, take the row sum formed from the bit of $a$ at $q$ and $b$, at indices
$j,p+q$, treating a missing bit as $0$. The sum of these row sums is at most the length
of $a$.
\end{lemma}

\begin{lemma}[Column equals the sum of rows]
\lean{Complexity.multRows_column}
\label{Complexity.multRows_column}
\leanok
\uses{Complexity.multColumn, Complexity.multRow, Complexity.multRows, Complexity.mult_column_double}
\difficulty{1}
Let $a,b$ be finite bit strings and let $j$ be a nonnegative integer. The sum, over
positions $q$ in $a$, of the row sum formed from the bit of $a$ at $q$ and $b$ at
indices $j,q$ equals the column sum formed from $a,b$ at position $j$. In both sums,
missing bits are taken to be $0$.
\end{lemma}

\begin{lemma}[Canonicality and the last bit]
\lean{Complexity.mult_canonical_last}
\label{Complexity.mult_canonical_last}
\leanok
\uses{Complexity.LogProg.Canon}
\difficulty{4}
A bit string $w$ is canonical exactly when it is empty or its entry at position $|w|-1$
is true, interpreting a missing entry as false.
\end{lemma}

\begin{lemma}[Characterization of multiplication encodings]
\lean{Complexity.mult_membership}
\label{Complexity.mult_membership}
\leanok
\uses{Complexity.LogProg.Canon, Complexity.LogProg.bitsVal, Complexity.LogProg.bitsVal_bits, Complexity.LogProg.canon_bits, Complexity.LogProg.canon_eq_bits, Complexity.multLang, Turing.pairEncode, Turing.pairEncode_injective}
\difficulty{4}
For any bit strings $a,b,c$, the nested pair encoding $\langle a,\langle
b,c\rangle\rangle$ belongs to the language of encoded products if and only if each of
$a,b,c$ is empty or ends in $1$, and the little-endian binary values of these words
satisfy $\operatorname{val}(c)=\operatorname{val}(a)\operatorname{val}(b)$, where
$\operatorname{val}(w_0\cdots w_{n-1})=\sum_{i=0}^{n-1}w_i2^i$.
\end{lemma}

\begin{lemma}[Characterisation of multiplication inputs]
\lean{Complexity.mult_acceptance}
\label{Complexity.mult_acceptance}
\leanok
\uses{Complexity.LogProg.Canon, Complexity.LogProg.canon_eq_bits, Complexity.multCarry, Complexity.multColumn, Complexity.multLang, Complexity.mult_columns_correct, Complexity.mult_membership, Complexity.mult_product_length, Turing.pairEncode}
\difficulty{4}
For bit strings $a,b,c$, the nested encoding $\langle a,\langle b,c\rangle\rangle$
belongs to the language of binary multiplication exactly when $a,b,c$ are canonical,
$|c|\le |a|+|b|$, the carry after column $|a|+|b|$ is zero, and for every $j<|a|+|b|$,
the $j$th bit of $c$ equals the remainder modulo $2$ of the sum of the carry into column
$j$ and the product-column value at $j$.
\end{lemma}

\begin{definition}[Multiplication machine states]
\lean{Complexity.MultState}
\label{Complexity.MultState}
\leanok
Define a finite state set with the following states: length testing and incrementing;
pair testing and first-bit handling, each with a Boolean inner flag; pair second-bit and
pair-reading states, each with Boolean inner and bit flags; pair advancement with a
Boolean inner flag and a choice in a four-element set; copying with a stage in an
eleven-element set and a Boolean bump flag; end-down states with Boolean inner and
second flags; half-counting and count-down states with a Boolean second flag;
canonical-form testing and clearing with a field in a three-element set, canonical down
states with a field and Boolean second flag, and canonical reading with a field; and the
parameter-free states column, clear-A, clear-P, outer, clear-B, clear-diagonal, inner,
diagonal, first-bit, second-bit, add, next-P, result-test, result-read, carry-half,
next-column, next-result-test, next-result, and final-position. Also include next-B with
a Boolean second flag, next-diagonal, next-A with a Boolean second flag, parity with a
Boolean bit, final-carry, and answer with a Boolean bit.
\end{definition}

\begin{definition}[Pair-processing state transition]
\lean{Complexity.multPairNext}
\label{Complexity.multPairNext}
\leanok
\uses{Complexity.MultState}
For an inner-mode flag $b$ and a choice $j\in\{0,1,2,3\}$, define the next state by the
table: choices $0$ and $1$ lead to the pair-second state with bit $0$ and bit $1$,
respectively; choice $2$ leads to the pair-test state; choice $3$ leads to copy stage
$1$ when $b$ is true and stage $0$ when $b$ is false, with the bump flag false.
\end{definition}

\begin{definition}[Multiplication copy-stage table]
\lean{Complexity.multCopy}
\label{Complexity.multCopy}
\leanok
\uses{Complexity.MultState}
Define a table on the eleven copy stages, with each entry consisting of two register
indices and a state. In order of stages $0$ through $10$, the entries are
$(2,0,\text{pair-test with inner flag true})$, $(3,0,\text{copy stage }2\text{, bump
false})$, $(4,2,\text{end-down with flags false, false})$, $(5,3,\text{end-down with
flags true, false})$, $(6,3,\text{count-half})$, $(0,4,\text{canonical-down field
}0\text{, second false})$, $(0,5,\text{canonical-down field }1\text{, second false})$,
$(0,1,\text{canonical-down field }2\text{, second false})$, $(12,3,\text{column})$,
$(9,2,\text{clear-diagonal})$, and $(10,8,\text{inner})$.
\end{definition}

\begin{definition}[Multiplication register program]
\lean{Complexity.multARM}
\label{Complexity.multARM}
\leanok
\uses{Complexity.LogProg.ARM, Complexity.MultState, Complexity.multCopy, Complexity.multPairNext}
Define a register program with $13$ registers, two oracle deciders, and the finite set
of states specified for the multiplication procedure. In each state, its instruction is
given by the displayed transition table: the program uses register increments,
decrements, clearing, halving, zero, parity, and equality tests, and calls the indicated
oracle on the specified virtual input; its terminal states return the indicated Boolean
answer.
\end{definition}

\begin{lemma}[Addresses in a nested encoding]
\lean{Complexity.mult_addresses}
\label{Complexity.mult_addresses}
\leanok
\uses{Complexity.pairEncode_getD, Complexity.pairEncode_getD_snd, Turing.pairEncode}
\difficulty{4}
For bit strings $a,b,c$, the nested encoding $\langle a,\langle b,c\rangle\rangle$ has
the following entries: at position $2i$ it equals the $i$th bit of $a$ whenever $i<|a|$;
at position $2|a|+2+2i$ it equals the $i$th bit of $b$ whenever $i<|b|$; and at position
$2|a|+2+2|b|+2+i$ it equals the $i$th bit of $c$ for every natural number $i$, with
out-of-range list entries interpreted as false.
\end{lemma}

\begin{definition}[Remaining multiplication bits]
\lean{Complexity.MultRemaining}
\label{Complexity.MultRemaining}
\leanok
\uses{Complexity.multCarry, Complexity.multColumn}
For Boolean lists $a,b,c$ and a natural number $j$, define the remaining-bits condition
to mean that $c$ has length at most $|a|+|b|$, the carry after position $|a|+|b|$ is
zero, and at every position $i$ with $j\le i<|a|+|b|$, the bit of $c$ (taking a missing
bit to be false) equals the sum of the carry and column value at position $i$, reduced
modulo $2$.
\end{definition}

\begin{lemma}[One-step characterization of remaining product bits]
\lean{Complexity.mult_remaining_step}
\label{Complexity.mult_remaining_step}
\leanok
\uses{Complexity.MultRemaining, Complexity.multCarry, Complexity.multColumn}
\difficulty{4}
Let $a,b,c$ be finite bit strings and let $j<|a|+|b|$. The following are equivalent:
$|c|\le |a|+|b|$, the carry after $|a|+|b|$ columns is zero, and every position $i$ from
$j$ through $|a|+|b|-1$ has bit of $c$ equal to the parity of the carry into column $i$
plus the sum of the bit products contributing to column $i$; or, the bit of $c$ at
position $j$ equals that column’s parity and the same remaining-bit conditions hold from
position $j+1$ onward.
\end{lemma}

\begin{definition}[Multiplication oracle deciders]
\lean{Complexity.multOracle}
\label{Complexity.multOracle}
\leanok
\uses{Complexity.inputBitLang, Complexity.inputLenLang, Turing.MultiTapeTM.indicator}
For each of the two deciders, define its Boolean answer on a Boolean-list input to be
the indicator of the corresponding input language: the first decider uses the language
of valid multiplication inputs, and the second uses the language of valid input lengths.
\end{definition}

\begin{lemma}[Bit and length oracle queries]
\lean{Complexity.mult_queries}
\label{Complexity.mult_queries}
\leanok
\uses{Complexity.indicator_indexLang_pairEncode, Complexity.inputBitLang, Complexity.inputLenLang, Complexity.multOracle, Turing.pairEncode}
\difficulty{1}
For every finite bit string $x$ and natural number $i$, querying the first oracle on the
encoded pair consisting of $x$ and the binary word for $i$ returns the bit of $x$ at
position $i$, with false returned when that position is out of range. Querying the
second oracle on the same encoded pair returns true exactly when $i<|x|$.
\end{lemma}

\begin{lemma}[Instruction side conditions]
\lean{Complexity.mult_preconditions}
\label{Complexity.mult_preconditions}
\leanok
\uses{Complexity.LogProg.PreS, Complexity.MultState, Complexity.multARM}
\difficulty{4}
For every bit string $x$, every state $s$ of the fixed multiplication program, every
assignment of natural numbers to its thirteen registers, and either choice of recorded
answer, the resulting configuration satisfies the program's instruction side conditions:
register comparisons use distinct registers, and any list of registers supplied to an
oracle call has no repetitions.
\end{lemma}

\begin{definition}[Good multiplication configurations]
\lean{Complexity.MultGood}
\label{Complexity.MultGood}
\leanok
\uses{Complexity.LogProg.AConf, Complexity.LogProg.PreS, Complexity.MultState, Complexity.multARM}
For a Boolean-list input $x$ and a configuration with $13$ registers and a
multiplication state, define the configuration to be good when it is a permitted initial
configuration for the multiplication program on $x$ and every register value is at most
$4(|x|+1)$.
\end{definition}

\begin{definition}[Reachable multiplication configurations]
\lean{Complexity.MultReach}
\label{Complexity.MultReach}
\leanok
\uses{Complexity.LogProg.AReach, Complexity.MultGood, Complexity.multARM, Complexity.multOracle}
For a Boolean-list input $x$, define the reachable multiplication configurations to be
the configurations reachable from the allowed initial configurations for $x$ under the
multiplication program, using its two oracle deciders and restricting configurations to
those satisfying the good-configuration condition.
\end{definition}

\begin{definition}[Halting multiplication configurations]
\lean{Complexity.MultHalt}
\label{Complexity.MultHalt}
\leanok
\uses{Complexity.LogProg.AHalt, Complexity.MultGood, Complexity.multARM, Complexity.multOracle}
For a Boolean-list input $x$, define the multiplication halting configurations to be the
allowed configurations at which the multiplication program halts, using its two oracle
deciders and the good-configuration condition.
\end{definition}

\begin{lemma}[Copying one register value]
\lean{Complexity.mult_copy}
\label{Complexity.mult_copy}
\leanok
\uses{Complexity.LogProg.registerCopy_loop, Complexity.MultReach, Complexity.multARM, Complexity.multCopy, Complexity.multOracle}
\difficulty{1}
Let $x$ be a finite bit string, let $k$ be an index from $0$ to $10$, and let $v$ be a
vector of thirteen natural-number registers, each at most $4(|x|+1)$. Suppose the
register selected as the destination by stage $k$ of the fixed multiplication register
program is no larger than the register selected as its source. Starting at stage $k$ of
the copy routine with registers $v$ and no answer, the program can reach the stage
specified for $k$, with the destination register replaced by the source register's
original value, all other registers unchanged, and no answer.
\end{lemma}

\begin{lemma}[Scanning a paired bit string]
\lean{Complexity.mult_scan}
\label{Complexity.mult_scan}
\leanok
\uses{Complexity.LogProg.alignedPair_loop, Complexity.MultReach, Complexity.MultState, Complexity.multARM, Complexity.multOracle, Complexity.multPairNext, Complexity.mult_queries, Turing.pairDecode}
\difficulty{4}
Let $x,w$ be finite bit strings, let $b$ be either true or false, and let $v$ assign a
natural-number counter to each of the program’s 13 registers. Suppose every counter is
at most $4(|x|+1)$, register $1$ equals register $0$ plus $|w|$, and for every $j<|w|$
the bit of $x$ at position $v(0)+j$ (taken to be false when out of range) equals the bit
of $w$ at position $j$ (also taken to be false when out of range). If the paired-string
parser fails on $w$, the program can reach its false-answer state after advancing
register $0$ by some $t\le |w|$, while maintaining its safety conditions at every
earlier step. If the parser returns $(a,\_)$, it can reach the copy state indexed by $b$
after advancing register $0$ by $2|a|+2$, again maintaining those conditions at every
earlier step.
\end{lemma}

\begin{lemma}[One diagonal contribution to multiplication]
\lean{Complexity.mult_contribution}
\label{Complexity.mult_contribution}
\leanok
\uses{Complexity.LogProg.AReach.refl, Complexity.LogProg.AReach.single, Complexity.LogProg.AReach.step, Complexity.LogProg.PreS, Complexity.LogProg.astep, Complexity.LogProg.vword_whole₁, Complexity.MultReach, Complexity.MultState, Complexity.multARM, Complexity.mult_preconditions, Complexity.mult_queries}
\difficulty{5}
Let $x$ be a finite bit string and let $v$ be a vector of thirteen natural-number
registers, each at most $4(|x|+1)$. In the fixed multiplication register program, a
configuration in the diagonal-processing state with registers $v$ and no answer can
reach a configuration in the next-$b$ state with registers identical except that
register 11 is increased by $v_{11}+[v_{10}=v_7]x_{v_0}x_{v_9}$, and still with no
answer.
\end{lemma}

\begin{lemma}[Inner multiplication loop]
\lean{Complexity.mult_inner_loop}
\label{Complexity.mult_inner_loop}
\leanok
\uses{Complexity.LogProg.AReach.single, Complexity.LogProg.AReach.step, Complexity.LogProg.AReach.trans, Complexity.LogProg.astep, Complexity.LogProg.register_update_le, Complexity.MultReach, Complexity.multARM, Complexity.multOracle, Complexity.multRow, Complexity.multRow_cons, Complexity.mult_contribution, Complexity.mult_preconditions}
\difficulty{6}
Let $x$ and $w$ be bit strings, let $b\in\{0,1\}$, and let $v$ assign natural-number
values to the thirteen registers. Suppose every register is at most $4(|x|+1)$,
$v(5)=v(9)+2|w|$, the bit of $x$ at position $v(0)$ (taken to be $0$ if out of range) is
$b$, and for each $q<|w|$ the bit at position $v(9)+2q$ equals the $q$th bit of $w$
(again using $0$ out of range). Also suppose $v(10)+|w|\le 4(|x|+1)$ and $v(11)+R\le
4(|x|+1)$, where $R=\sum_{q<|w|,\ v(10)+q=v(7)} b\,w_q$ and bits are identified with $0$
or $1$. Then the fixed multiplication program reaches the state preceding the next
outer-loop step, with registers $9$, $10$, and $11$ respectively increased by $2|w|$,
$|w|$, and $R$, all other registers unchanged, and no answer recorded. Every earlier
configuration on this run satisfies the program's instruction side conditions and
register bound.
\end{lemma}

\begin{lemma}[Starting an inner multiplication row]
\lean{Complexity.mult_row_start}
\label{Complexity.mult_row_start}
\leanok
\uses{Complexity.LogProg.AReach.step, Complexity.LogProg.AReach.trans, Complexity.LogProg.register_update_le, Complexity.MultReach, Complexity.multCopy, Complexity.mult_copy, Complexity.mult_preconditions}
\difficulty{4}
Let $x$ be a finite bit string and let $v$ assign natural numbers to $13$ registers,
each at most $4(|x|+1)$. Starting in the state that clears register $9$, with register
values $v$ and no answer recorded, the multiplication machine can reach its inner-loop
state with register $9$ equal to the initial value of register $2$ and register $10$
equal to the initial value of register $8$, while satisfying its prescribed validity and
register-bound conditions at every intermediate step.
\end{lemma}

\begin{lemma}[Outer multiplication loop]
\lean{Complexity.mult_outer_loop}
\label{Complexity.mult_outer_loop}
\leanok
\uses{Complexity.LogProg.AReach.single, Complexity.LogProg.AReach.step, Complexity.LogProg.AReach.trans, Complexity.LogProg.astep, Complexity.LogProg.register_update_le, Complexity.MultReach, Complexity.multARM, Complexity.multOracle, Complexity.multRow, Complexity.multRows, Complexity.multRows_cons, Complexity.mult_inner_loop, Complexity.mult_preconditions, Complexity.mult_row_start}
\difficulty{6}
Let $x,a,b$ be bit strings and let $v$ be a register assignment bounded by $4(|x|+1)$ in
every register. Suppose $v(4)=v(0)+2|a|$, $v(5)=v(2)+2|b|$, the input bits at positions
$v(0)+2p$ and $v(2)+2q$ reproduce respectively the bits of $a$ and $b$ for every $p<|a|$
and $q<|b|$ (with out-of-range bits treated as $0$), and $v(8)+|a|+|b|\le 4(|x|+1)$.
Define $S=\sum_{p<|a|}\sum_{q<|b|} [v(7)=v(8)+p+q]a_p b_q$, where bits are identified
with $0$ or $1$, and suppose $v(11)+S\le 4(|x|+1)$. Then there is a final register
assignment, still bounded by $4(|x|+1)$, in which registers $0$, $8$, and $11$ equal
respectively $v(0)+2|a|$, $v(8)+|a|$, and $v(11)+S$, and every register other than
$0,8,9,10,11$ retains its initial value. The fixed multiplication program reaches its
result-test state with this assignment and no answer recorded, with all earlier
configurations satisfying its instruction side conditions and register bound.
\end{lemma}

\begin{lemma}[Returning a designated answer]
\lean{Complexity.mult_return}
\label{Complexity.mult_return}
\leanok
\uses{Complexity.LogProg.AHalt.ret, Complexity.MultHalt, Complexity.mult_preconditions}
\difficulty{1}
Let $x$ be a finite bit string, let $b$ be a Boolean value, and let $v$ assign natural
numbers to $13$ registers, each at most $4(|x|+1)$. Starting in the answer state
labelled $b$ with register values $v$ and no answer recorded, the multiplication machine
halts with answer $b$.
\end{lemma}

\begin{lemma}[Termination at the final position]
\lean{Complexity.mult_finish}
\label{Complexity.mult_finish}
\leanok
\uses{Complexity.LogProg.AHalt.step, Complexity.LogProg.astep, Complexity.MultHalt, Complexity.multARM, Complexity.mult_preconditions, Complexity.mult_return}
\difficulty{4}
For every bit string $x$ and register assignment $v$ with $v(r)\le 4(|x|+1)$ for every
register $r$, the fixed multiplication program, started in its final-position state with
registers $v$ and no answer recorded, halts with Boolean answer true exactly when
$v(12)=v(1)$ and $v(11)=0$. Every configuration before halting satisfies the program's
instruction side conditions and the same register bound.
\end{lemma}

\begin{lemma}[Parity step in multiplication]
\lean{Complexity.mult_parity}
\label{Complexity.mult_parity}
\leanok
\uses{Complexity.LogProg.AHalt.step, Complexity.LogProg.AReach.single, Complexity.LogProg.astep, Complexity.MultHalt, Complexity.MultReach, Complexity.multARM, Complexity.multOracle, Complexity.mult_preconditions, Complexity.mult_return}
\difficulty{5}
Let $x$ be a bit string, $v$ a register assignment bounded by $4(|x|+1)$ in every
register, and $b\in\{0,1\}$. If $b\equiv v(11)\pmod 2$, the fixed multiplication program
reaches the carry-halving state from the parity state with bit $b$, without changing the
registers or recording an answer. Otherwise it halts with answer false.
\end{lemma}

\begin{lemma}[Reading a product result bit]
\lean{Complexity.mult_result_read}
\label{Complexity.mult_result_read}
\leanok
\uses{Complexity.LogProg.AReach.refl, Complexity.LogProg.AReach.single, Complexity.LogProg.AReach.step, Complexity.LogProg.PreS, Complexity.LogProg.astep, Complexity.LogProg.vword_whole₁, Complexity.MultReach, Complexity.multARM, Complexity.multOracle, Complexity.mult_preconditions, Complexity.mult_queries}
\difficulty{4}
Let $x$ be a finite bit string and let $v$ assign natural numbers to $13$ registers,
with every register value at most $4(|x|+1)$ and register $1$ equal to $|x|$. Starting
in the result-test state with register values $v$ and no answer recorded, the
multiplication machine can reach the parity state labelled by the bit of $x$ at position
$v(12)$, while satisfying its prescribed validity and register-bound conditions at every
intermediate step.
\end{lemma}

\begin{lemma}[Advance to the next column]
\lean{Complexity.mult_next_column}
\label{Complexity.mult_next_column}
\leanok
\uses{Complexity.LogProg.AReach.refl, Complexity.LogProg.AReach.single, Complexity.LogProg.AReach.step, Complexity.LogProg.astep, Complexity.LogProg.register_update_le, Complexity.MultReach, Complexity.multARM, Complexity.multOracle, Complexity.mult_preconditions}
\difficulty{5}
Let $x$ be a bit string and $v$ a register assignment with every register at most
$4(|x|+1)$. If $v(12)\le v(1)$ and $v(7)+1\le 4(|x|+1)$, then from the carry-halving
state the fixed multiplication program reaches the column state with registers $11$,
$7$, and $12$ replaced respectively by $\lfloor v(11)/2\rfloor$, $v(7)+1$, and
$\min(v(12)+1,v(1))$; all other registers are unchanged, and no answer is recorded.
Every earlier configuration satisfies the program's instruction side conditions and
register bound.
\end{lemma}

\begin{lemma}[Execution of the multiplication procedure]
\lean{Complexity.mult_columns_run}
\label{Complexity.mult_columns_run}
\leanok
\uses{Complexity.LogProg.AHalt.step, Complexity.LogProg.AReach.halt, Complexity.LogProg.AReach.single, Complexity.LogProg.AReach.step, Complexity.LogProg.astep, Complexity.LogProg.register_update_le, Complexity.MultHalt, Complexity.MultReach, Complexity.MultRemaining, Complexity.multARM, Complexity.multCarry, Complexity.multCarry_le, Complexity.multColumn, Complexity.multColumn_le, Complexity.multRows_column, Complexity.mult_addresses, Complexity.mult_finish, Complexity.mult_next_column, Complexity.mult_outer_loop, Complexity.mult_parity, Complexity.mult_preconditions, Complexity.mult_remaining_step, Complexity.mult_result_read, Turing.length_pairEncode, Turing.pairEncode}
\difficulty{7}
Let $a,b,c,x$ be finite bit strings with $x=\langle a,\langle b,c\rangle\rangle$, and
let $j\le |a|+|b|$. Let $v$ be a vector of thirteen natural-number registers, each at
most $4(|x|+1)$, satisfying $v_1=|x|$, $v_2=2|a|+2$, $v_4=2|a|$, $v_5=2|a|+2+2|b|$, and
$v_6=|a|+|b|$. Suppose also that $v_7=j$, $v_{11}=K_j$, and
$v_{12}=2|a|+2+2|b|+2+\min(j,|c|)$, where $K_j$ is the carry into column $j$ for
multiplying $a$ and $b$. The fixed multiplication procedure, started in its
column-processing state with registers $v$ and no answer yet, halts with the Boolean
value of the following condition: $|c|\le |a|+|b|$, $K_{|a|+|b|}=0$, and for every $i$
with $j\le i<|a|+|b|$, the $i$th bit of $c$ equals $(K_i+C_i(a,b))$ modulo $2$.
\end{lemma}

\begin{definition}[Canonical source positions]
\lean{Complexity.multCanonSource}
\label{Complexity.multCanonSource}
\leanok
Define a map from $\mathrm{Fin}(3)$ to $\mathrm{Fin}(13)$ by sending $0$ to $4$, $1$ to $5$, and
$2$ to $1$.
\end{definition}

\begin{definition}[Canonical copy positions]
\lean{Complexity.multCanonCopy}
\label{Complexity.multCanonCopy}
\leanok
Define $c:\mathrm{Fin}(3)\to\mathrm{Fin}(11)$ by $c(0)=5$, $c(1)=6$, and $c(2)=7$.
\end{definition}

\begin{definition}[Next canonical multiplication state]
\lean{Complexity.multCanonNext}
\label{Complexity.multCanonNext}
\leanok
\uses{Complexity.MultState}
For $f\in\{0,1,2\}$, define a state $q(f)$ of the multiplication procedure as follows:
if $f=0$, it is the canonical test state for field $1$; if $f=1$, it is the canonical
test state for field $2$; and if $f=2$, it is the copy state at stage $8$ with its bump
flag false.
\end{definition}

\begin{lemma}[Final-bit processing]
\lean{Complexity.mult_last_bit}
\label{Complexity.mult_last_bit}
\leanok
\uses{Complexity.LogProg.AReach.single, Complexity.LogProg.AReach.step, Complexity.LogProg.AReach.trans, Complexity.LogProg.astep, Complexity.LogProg.register_update_le, Complexity.LogProg.vword_whole₁, Complexity.MultReach, Complexity.multARM, Complexity.multCanonCopy, Complexity.multCanonNext, Complexity.multCanonSource, Complexity.multCopy, Complexity.multOracle, Complexity.mult_copy, Complexity.mult_preconditions, Complexity.mult_queries}
\difficulty{4}
For every bit string $x$, every $f\in\{0,1,2\}$, and every register assignment $v$
bounded by $4(|x|+1)$, the fixed multiplication program, started in the clearing state
for field $f$, reaches a configuration with register $0$ set to $v(s_f)-c_f$, where
$(s_0,s_1,s_2)=(4,5,1)$ and $(c_0,c_1,c_2)=(2,2,1)$. If the bit of $x$ at that position
is $1$ (with out-of-range bits treated as $0$), the resulting state is the next
canonical-field test; otherwise it is the rejecting answer state. Every earlier
configuration satisfies the program's instruction side conditions and register bound.
\end{lemma}

\begin{lemma}[Validation of a canonical field]
\lean{Complexity.mult_validate_field}
\label{Complexity.mult_validate_field}
\leanok
\uses{Complexity.LogProg.AHalt.step, Complexity.LogProg.AReach.halt, Complexity.LogProg.Canon, Complexity.LogProg.astep, Complexity.LogProg.register_update_le, Complexity.MultHalt, Complexity.multARM, Complexity.multCanonNext, Complexity.multCanonSource, Complexity.multOracle, Complexity.mult_canonical_last, Complexity.mult_last_bit, Complexity.mult_preconditions, Complexity.mult_return}
\difficulty{5}
Let $x,w$ be finite bit strings, let $f$ be one of the three field indices, let $v$ be
an assignment of natural numbers to the 13 registers, and let $b$ be a Boolean answer.
Assume every register is at most $4(|x|+1)$. Suppose that one program step from the
field-test state goes to the next validation state when $w$ is empty, and otherwise goes
to the field-clear state. If $w$ is nonempty, suppose the bit of $x$ at the position
obtained from the designated source register by subtracting $1$ when $f$ is the third
field and $2$ otherwise (with natural-number subtraction truncated at zero) agrees with
the last bit of $w$, taking out-of-range bits to be false. Assume also that a
noncanonical $w$ forces $b$ to be false, and that for canonical $w$, every bounded
register assignment agreeing with $v$ on all registers other than register $0$ yields a
halting run from the next validation state with answer $b$. Then the run from the
field-test state halts with answer $b$, maintaining the program’s safety conditions
before halting.
\end{lemma}

\begin{definition}[Ready multiplication configuration]
\lean{Complexity.MultReady}
\label{Complexity.MultReady}
\leanok
Given Boolean lists $x,a,b,c$ and a function $v$ assigning a natural-number value to
each position in $\{0,1,\ldots,12\}$, the configuration is ready when $v(1)=|x|$,
$v(2)=2|a|+2$, $v(3)=2|a|+2+2|b|+2$, $v(4)=2|a|$, $v(5)=2|a|+2+2|b|$, $v(6)=|a|+|b|$,
and $v(7)=v(11)=v(12)=0$.
\end{definition}

\begin{lemma}[Preservation of the multiplication frame]
\lean{Complexity.mult_ready_frame}
\label{Complexity.mult_ready_frame}
\leanok
\uses{Complexity.MultReady}
\difficulty{1}
Let $x,a,b,c$ be finite Boolean words, and let $v,u:\{0,1,\ldots,12\}\to\mathbb{N}$. If
$v$ satisfies the ready multiplication conditions for $x,a,b,c$ and $u(r)=v(r)$ for
every $r\ne 0$, then $u$ also satisfies those conditions.
\end{lemma}

\begin{lemma}[Correctness of the Multiplication Computation]
\lean{Complexity.mult_canonical_run}
\label{Complexity.mult_canonical_run}
\leanok
\uses{Complexity.LogProg.AReach.halt, Complexity.LogProg.Canon, Complexity.LogProg.astep, Complexity.LogProg.register_update_le, Complexity.MultHalt, Complexity.MultReady, Complexity.MultRemaining, Complexity.multARM, Complexity.multCanonNext, Complexity.multCarry, Complexity.multCopy, Complexity.multLang, Complexity.mult_acceptance, Complexity.mult_addresses, Complexity.mult_columns_run, Complexity.mult_copy, Complexity.mult_ready_frame, Complexity.mult_validate_field, Turing.length_pairEncode, Turing.pairEncode}
\difficulty{6}
Let $x,a,b,c$ be Boolean words such that $x=\langle a,\langle b,c\rangle\rangle$, and
let $v$ assign a natural number to each of thirteen registers. Suppose every register
value is at most $4(|x|+1)$, and the register values satisfy $v(1)=|x|$, $v(2)=2|a|+2$,
$v(3)=2|a|+2+2|b|+2$, $v(4)=2|a|$, $v(5)=2|a|+2+2|b|$, $v(6)=|a|+|b|$, and
$v(7)=v(11)=v(12)=0$. Starting in the canonicality-test state for field $0$, with
register values $v$ and no answer recorded, the multiplication procedure eventually
halts with the Boolean answer indicating whether $x$ belongs to the multiplication
language.
\end{lemma}

\begin{lemma}[Two decrements in succession]
\lean{Complexity.mult_sub_two}
\label{Complexity.mult_sub_two}
\leanok
\uses{Complexity.LogProg.AReach.single, Complexity.LogProg.AReach.step, Complexity.LogProg.astep, Complexity.LogProg.register_update_le, Complexity.MultReach, Complexity.MultState, Complexity.multARM, Complexity.multOracle, Complexity.mult_preconditions}
\difficulty{3}
Let $x$ be a finite bit string, let $v$ assign a natural number to each of the 13
registers, let $r$ be a register, and let $\ell_0,\ell_1$ be two program states, with a
destination state $d$. Suppose the program instruction at $\ell_0$ decrements register
$r$ and moves to $\ell_1$, and the instruction at $\ell_1$ decrements register $r$ and
moves to $d$. If every register initially has value at most $4(|x|+1)$, then the program
can reach state $d$ from state $\ell_0$ with register $r$ equal to its initial value
truncated after subtracting two, all other registers unchanged, and the safety
conditions maintained at every earlier step.
\end{lemma}

\begin{lemma}[Multiplication machine halting]
\lean{Complexity.mult_prepare_run}
\label{Complexity.mult_prepare_run}
\leanok
\uses{Complexity.LogProg.AReach.halt, Complexity.LogProg.AReach.single, Complexity.LogProg.astep, Complexity.LogProg.register_update_le, Complexity.MultHalt, Complexity.MultReach, Complexity.MultReady, Complexity.multARM, Complexity.multCopy, Complexity.multLang, Complexity.multOracle, Complexity.mult_canonical_run, Complexity.mult_copy, Complexity.mult_preconditions, Complexity.mult_sub_two, Turing.pairEncode}
\difficulty{5}
Let $x,a,b,c$ be finite bit strings with $x=\langle a,\langle b,c\rangle\rangle$, and
let $v$ assign a natural number to each of the machine’s $13$ registers. Suppose every
register value is at most $4(|x|+1)$, register $1$ contains $|x|$, registers $2$ and $3$
contain $2|a|+2$ and $2|a|+2+2|b|+2$, respectively, and registers $4,5,6,7,11,12$ are
all zero. Starting in the copy stage with these register values and no answer recorded,
the multiplication machine halts with answer true exactly when $x$ belongs to the
language of encoded multiplication instances.
\end{lemma}

\begin{lemma}[Failure to decode excludes multiplication inputs]
\lean{Complexity.mult_bad_decode}
\label{Complexity.mult_bad_decode}
\leanok
\uses{Complexity.multLang, Turing.pairDecode, Turing.pairDecode_pairEncode}
\difficulty{4}
For a bit string $x$, if parsing $x$ as a pair and then parsing the second component
fails, then $x$ does not belong to the language of binary multiplication.
\end{lemma}

\begin{lemma}[Multiplication program correctness]
\lean{Complexity.mult_run}
\label{Complexity.mult_run}
\leanok
\uses{Complexity.LogProg.AReach.halt, Complexity.LogProg.inputLength_loop, Complexity.LogProg.register_update_le, Complexity.MultHalt, Complexity.MultReach, Complexity.multARM, Complexity.multCopy, Complexity.multLang, Complexity.multOracle, Complexity.mult_bad_decode, Complexity.mult_copy, Complexity.mult_prepare_run, Complexity.mult_queries, Complexity.mult_return, Complexity.mult_scan, Complexity.pairEncode_getD, Turing.eq_pairEncode_of_pairDecode, Turing.length_pairEncode, Turing.pairDecode, Turing.pairEncode}
\difficulty{6}
For every finite bit string $x$, the specified 13-register oracle program, started in
its length-testing state with all registers zero and no answer recorded, halts with
answer true exactly when $x$ encodes a pair of binary numerals whose product is the
third numeral. Throughout every step before halting, its configuration satisfies the
program’s safety conditions and every register is at most $4(|x|+1)$.
\end{lemma}

\begin{theorem}[Multiplication language is in log space]
\lean{Complexity.multLang_mem_LOGSPACE}
\label{Complexity.multLang_mem_LOGSPACE}
\leanok
\uses{Complexity.LOGSPACE, Complexity.LogProg.arm_decides_poly, Complexity.LogProg.indicator_eq_decide, Complexity.inputBitLang, Complexity.inputBitLang_mem, Complexity.inputLenLang, Complexity.inputLenLang_mem, Complexity.multARM, Complexity.multLang, Complexity.multOracle, Complexity.mult_run}
\difficulty{4}
Let $L$ be the language of bit strings of the form $\langle
\operatorname{bits}(a),\langle
\operatorname{bits}(b),\operatorname{bits}(ab)\rangle\rangle$, where $a,b$ are
nonnegative integers, $\operatorname{bits}(n)$ is the bit string of $n$, and $\langle
x,y\rangle$ doubles each bit of $x$, appends the separator $01$, and then appends $y$.
Then $L$ is decidable by a finite-state, deterministic multi-tape Turing machine using
at most $c(\lfloor\log_2 n\rfloor+1)$ work-tape cells on inputs of length $n$, for some
nonnegative integer constant $c$.
\end{theorem}
